\section{Troubleshooting protocols}

The protocol editor's \textbf{Protocol checks} panel separates errors from
warnings. Start with the first error because one malformed instruction can
make later messages less useful.

\subsection{Unknown instruction}

\textbf{Problem fragment}
\begin{lstlisting}
FLIP -I Q -O Q
\end{lstlisting}

\textbf{Typical error:} \code{Unknown command: FLIP}

\textbf{Why:} FLIP is not a QPU instruction. The intended bit-flip gate is X.

\textbf{Corrected fragment}
\begin{lstlisting}
X -I Q -O Q
\end{lstlisting}

\subsection{Missing gate output}

\textbf{Problem fragment}
\begin{lstlisting}
H -I Q
\end{lstlisting}

\textbf{Typical error:} H requires an \code{-O} output.

\textbf{Why:} a one-qubit gate names the same flowing wire as both input and
mutated output.

\textbf{Corrected fragment}
\begin{lstlisting}
H -I Q -O Q
\end{lstlisting}

\subsection{Wrong number of controls}

\textbf{Problem fragment}
\begin{lstlisting}
CNOT -I A B -O Target
\end{lstlisting}

\textbf{Why:} CNOT has one control. A two-control bit flip is CCNOT.

\textbf{Choose the correction that matches the intent}
\begin{lstlisting}
CNOT  -I A   -O Target  # one control
CCNOT -I A B -O Target  # two controls
\end{lstlisting}

\subsection{Invalid phase angle}

\textbf{Problem fragment}
\begin{lstlisting}
PHASE=ninety -I Q -O Q
\end{lstlisting}

\textbf{Why:} the angle must be numeric radians, degrees ending in d, or pi
notation.

\textbf{Equivalent corrected fragments}
\begin{lstlisting}
PHASE=pi/2   -I Q -O Q
PHASE=90d    -I Q -O Q
PHASE=1.5708 -I Q -O Q
\end{lstlisting}

\subsection{Unknown child process}

\textbf{Problem fragment}
\begin{lstlisting}
RUNCHILD Adder -I A B Cin -O Carry Sum
\end{lstlisting}

\textbf{Typical error:} \code{Unknown child process 'Adder'}

\textbf{Why:} child source must exist in the process catalog and its
MAIN-PROCESS name must match exactly. Load the child protocol first, or change
Adder to the catalog's declared name.

\subsection{Truth-table columns do not match}

If protocol source declares
\begin{lstlisting}
PARAMS: A:state B:state
...
RETURNVALS Carry Sum
\end{lstlisting}
then companion metadata must keep exactly that input and output order:
\begin{lstlisting}
INPUTS: A B
OUTPUTS: Carry Sum
# A B Carry Sum
\end{lstlisting}
Renaming, reordering, or omitting a column changes the interface and causes
validation errors.

\subsection{Cycle-suffix warnings}

An integer token suffix is checked against the process cycle.
\code{Q:1} while the process is still on cycle 0 produces
\code{CYCLE_SUFFIX_MISMATCH}. Child processes count cycles from zero again.
\code{-$R}, JOIN, SPLIT, FREE, DELETETOKEN, COMPILEPROCESS, MASTERVAL,
\code{SAVE_STATE}, and \code{LOAD_STATE} now execute; they are not
metadata-only warnings. \code{DECLARECHILD} still binds a catalog child for
later \code{RUNCHILD} or \code{CALL}.

\subsection{Using the Correction Lab}

\begin{enumerate}
  \item In the protocol editor, read the reported line and suggestion.
  \item Choose \textbf{Open Correction Lab for guided help}.
  \item Confirm the selected process and its truth table. For a new experiment,
  use an uploaded or newly created process rather than modifying a protected
  bundled table.
  \item Run tests or probe outputs to identify the first behavioral mismatch.
  \item Apply gate-level guidance to protocol source.
  \item Return to the Circuit builder, compile again, and confirm that Protocol
  checks no longer reports an error.
\end{enumerate}

The Correction Lab is appropriate for syntax failures, child-process binding
problems, and circuits whose measured outputs disagree with their truth table.
It must not rewrite canonical bundled truth-table expectations.

\section{Diagnostic message reference}

\begin{longtable}{>{\ttfamily}p{0.25\textwidth}p{0.25\textwidth}p{0.4\textwidth}}
\toprule
\normalfont\bfseries Code & \bfseries Meaning & \bfseries First response\\
\midrule
\endhead
EMPTY\_PROTOCOL & No instructions were found. & Add a complete protocol
starting with MAIN-PROCESS.\\
EMPTY\_PARAMS & PARAMS has no name:type entries. & Add an entry such as
Input:state or remove PARAMS.\\
MISPLACED\_PARAMS & PARAMS is not the first instruction. & Move it before
MAIN-PROCESS.\\
UNKNOWN\_INSTRUCTION & The first word is not supported. & Check spelling and
the language reference.\\
INVALID\_PHASE & PHASE has an unreadable angle. & Use pi/2, 90d, or numeric
radians.\\
INVALID\_GATE\_ARITY & Input/control or output count is wrong. & Compare the
gate's documented \code{-I}/\code{-O} counts.\\
COMPILE\_ERROR & Syntax was readable but compilation could not complete. &
Inspect child names, SET values, and SWAP bindings.\\
CYCLE\_SUFFIX\_MISMATCH & An integer \code{:} suffix does not match the
process cycle. & Move the reference or add \code{INCREASECYCLE}.\\
INACTIVE\_REVERSE\_PREFIX & A B-prefixed gate other than S, T, or PHASE is
not a distinct inverse. & Drop the prefix on a self-inverse gate.\\
MISSING\_MAIN\_PROCESS & InlineProcess fallback will be used. & Add a stable
MAIN-PROCESS name.\\
MISSING\_RETURN\_VALUES & No public outputs are declared. & Add RETURNVALS in
display order.\\
\bottomrule
\end{longtable}
